\documentclass{article}
\usepackage[utf8]{inputenc}
\usepackage{a4wide}

\begin{document}

\section*{สี่เหลี่ยมใหญ่ที่สุดใน Histogram}

กำหนดอาร์เรย์ของจำนวนเต็มแทนความสูงของแท่งฮิสโตแกรม\textbf{โดยที่ความกว้างของแต่ละแท่งคือ

1}ให้\ul{\textbf{แสดงพื้นที่ของสี่เหลี่ยมผืนผ้าที่ใหญ่ที่สุดในฮิสโตแกรม}}เขียนโปรแกรมเพื่อรับข้อมูลเข้าทาง

stdin แล้วส่งไปเรียกฟังก์ชันที่ทำงานดังกล่าว คืนค่ากลับไปให้ main() แสดงผลลัพธ์\\



\hfill\break



{ตัวอย่างที่ 1}



\begin{quote}

\textbf{Input:}



\textbf{4}



\textbf{4 3 2 4}



\textbf{Output:~8}



\textbf{\includegraphics[width=2.73958in,height=2.08333in,alt={Screenshot 2567-03-19 at 23.45.38.png}]{/public/upload/e421d6c570.png}\\

}

\end{quote}



{{ตัวอย่างที่ 2}\\

}



\begin{quote}

\textbf{Input:}



\textbf{10}



\textbf{4 3 5 6 1 3 7 4 6 4}



\textbf{Output: 16}



\textbf{\includegraphics[width=2.60417in,height=2.08333in,alt={Screenshot 2567-03-19 at 23.46.29.png}]{/public/upload/e6ece751b7.png}\\

}

\end{quote}



\subsection*{Input}
บรรทัดแรกคือจำนวน{แท่งฮิสโตแกรม}



{บรรทัดสองคือ}ตัวเลขแสดงความสูงของ{แท่งฮิสโตแกรม}แต่ละตัวคั่นด้วยช่องว่าง



\subsection*{Output}
พื้นที่ของสี่เหลี่ยมผืนผ้าที่ใหญ่ที่สุดในฮิสโตแกรม\\



\end{document}
